% ══════════════════════════════════════════════════════════════════════════════
\section{Introduction}
% ══════════════════════════════════════════════════════════════════════════════

Ce rapport présente les résultats expérimentaux obtenus lors de la comparaison de trois modèles
d'exécution pour applications web Java~: Spring MVC avec JPA (bloquant classique), Spring WebFlux
avec R2DBC (réactif non bloquant) et Spring MVC avec Virtual Threads Java~21 (impératif non
bloquant).

L'expérimentation suit le protocole commun défini préalablement~\cite{protocole}~: même application
de gestion de commandes, même base PostgreSQL, même jeu de données (10~000 clients, 5~000 produits,
100~000 commandes) et mêmes scripts k6. Seul le modèle d'exécution varie entre les versions.

% ══════════════════════════════════════════════════════════════════════════════
\section{Environnement d'exécution}
% ══════════════════════════════════════════════════════════════════════════════

Les tests ont été exécutés sur une machine virtuelle commune conformément au protocole~:

\begin{center}
\begin{tabular}{|l|l|}
\hline
\textbf{Paramètre} & \textbf{Valeur} \\
\hline
Système                & Ubuntu Server 24.04 LTS \\
CPU                    & 4 vCPU \\
Mémoire                & 8 Go RAM \\
Java                   & Java 21 \\
Build                  & Maven \\
Tests de charge        & k6 \\
Base de données        & PostgreSQL (port 5432) \\
Service externe simulé & Mock HTTP (port 9090, latence contrôlée) \\
\hline
\end{tabular}
\end{center}

Les applications ont été exécutées séparément, une à la fois, avec les mêmes limites de ressources
(2~vCPU, 1,5~Go RAM par application). Les versions testées et leur configuration de pool de
connexions sont résumées ci-dessous~:

\begin{center}
\begin{tabular}{|L{3.5cm}|c|L{4.5cm}|c|c|}
\hline
\textbf{Version} & \textbf{Port} & \textbf{Technologies} & \textbf{Pool DB (utilisé)} & \textbf{Pool DB (cible)} \\
\hline
Spring MVC + JPA       & 8081 & Spring MVC, Hibernate, JDBC, HikariCP        & 20 & 50 \\
Spring WebFlux + R2DBC & 8082 & WebFlux, Project Reactor, R2DBC, WebClient   & 10\textsuperscript{*} & 50 \\
Spring MVC + VT        & 8083 & Spring MVC, JPA, Java~21 Virtual Threads     & 50 & 50 \\
\hline
\multicolumn{5}{|l|}{\small \textsuperscript{*}Valeur par défaut de \texttt{r2dbc-pool} (\texttt{spring.r2dbc.pool.max-size}~: 10). Aucune configuration explicite.} \\
\hline
\end{tabular}
\end{center}

\noindent\textbf{Déviation au protocole.}
Les trois versions ont été exécutées avec des tailles de pool différentes~: 20~connexions
(HikariCP) pour MVC JPA, 10~connexions (R2DBC pool par défaut) pour WebFlux et 50~connexions
(HikariCP) pour Virtual Threads. Le protocole prévoyait une configuration identique pour toutes
les versions. La valeur cible retenue pour les prochains passages est \textbf{50~connexions pour
chaque version}. Les résultats présentés dans ce rapport sont ceux des benchmarks tels qu'exécutés~;
les écarts liés à cette asymétrie sont discutés en section~\ref{sec:impl}.

% ══════════════════════════════════════════════════════════════════════════════
\section{Protocole de mesure}
% ══════════════════════════════════════════════════════════════════════════════

Quatre scénarios k6 ont été exécutés conformément au protocole commun~:

\begin{center}
\begin{tabular}{|l|L{4.5cm}|c|L{3.8cm}|}
\hline
\textbf{Scénario} & \textbf{Endpoints sollicités} & \textbf{VU} & \textbf{Objectif} \\
\hline
CRUD classique &
  GET /products/\{id\}, GET/POST /orders, GET /customers/\{id\}/orders &
  200 & Charge mixte réaliste \\
I/O-bound &
  GET /external/recommendations/\{id\} (300~ms simulée) &
  500 & Concurrence sous attente réseau \\
CPU-bound &
  GET /stress/cpu?n=35 (Fibonacci récursif) &
  50  & Facteur limitant~: CPU \\
Saturation &
  GET /products/\{id\} sans pause d'itération &
  800 & Débit maximal soutenable \\
\hline
\end{tabular}
\end{center}

Les métriques collectées sont~: débit (req/s), latence moyenne, médiane, P95, maximum et taux
d'erreur HTTP. Toutes les vérifications (\emph{checks}) ont été satisfaites à 100~\% pour
l'ensemble des scénarios et des versions.

% ══════════════════════════════════════════════════════════════════════════════
\section{Résultats expérimentaux}
% ══════════════════════════════════════════════════════════════════════════════

\subsection{Scénario CRUD classique (200~VU)}

Ce scénario reproduit une charge métier standard~: 60~\% de lectures produit, 20~\% de lectures
commande, 10~\% de lectures commandes client et 10~\% de créations de commande. L'itération intègre
une pause d'une seconde, ce qui plafonne le débit par la cadence du test et non par la capacité du
serveur.

\begin{center}
\begin{tabular}{|l|r|r|r|r|r|}
\hline
\textbf{Version} & \textbf{Req/s} & \textbf{Moy.~(ms)} & \textbf{Méd.~(ms)} & \textbf{P95~(ms)} & \textbf{Max~(ms)} \\
\hline
MVC JPA         & 106,78 &  3,51 &  2,51 &  10,12 &  477 \\
WebFlux         & 106,72 &  4,35 &  2,97 &  13,17 &   97 \\
Virtual Threads & 106,67 &  4,99 &  3,71 &  14,39 &   88 \\
\hline
\end{tabular}
\end{center}

Les trois versions atteignent un débit identique ($\approx$~106~req/s), confirmant que le goulot
est la cadence d'itération et non la capacité serveur. MVC JPA affiche la latence la plus faible
(P95 = 10~ms) mais présente un maximum de 477~ms révélant une pointe ponctuelle (saturation
momentanée du pool à 20~connexions ou pause GC). WebFlux et Virtual Threads montrent un maximum
nettement plus stable (97 et 88~ms respectivement).

\subsection{Scénario I/O-bound (500~VU, latence 300~ms)}

Ce scénario évalue la capacité à gérer la concurrence sous attente réseau prolongée. Chaque requête
attend 300~ms de réponse du service externe simulé, favorisant les architectures non bloquantes.

\begin{center}
\begin{tabular}{|l|r|r|r|r|r|}
\hline
\textbf{Version} & \textbf{Req/s} & \textbf{Moy.~(ms)} & \textbf{Méd.~(ms)} & \textbf{P95~(ms)} & \textbf{Max~(ms)} \\
\hline
MVC JPA         &  364,13 & 417,60 & 338,65 & 636,11 &   736 \\
WebFlux         &  444,13 & 306,24 & 303,41 & 320,47 &   396 \\
Virtual Threads &  445,33 & 304,69 & 302,44 & 314,82 &   423 \\
\hline
\end{tabular}
\end{center}

WebFlux et Virtual Threads sont statistiquement équivalents ($\approx$~444--445~req/s) et
maintiennent un P95 à seulement 15--20~ms au-dessus de la latence de base (300~ms). MVC JPA perd
18~\% de débit~: le pool de threads Tomcat s'épuise lorsque 500~VU bloquent simultanément sur
l'attente I/O, portant le P95 à 636~ms. Ce résultat est indépendant de la taille du pool de
connexions~: le goulot est ici le nombre de threads OS disponibles, non les connexions DB.

\subsection{Scénario CPU-bound (50~VU, Fibonacci n=35)}

Ce scénario soumet le serveur à un calcul intensif (Fibonacci récursif naïf, $n = 35$, sans
mémoïsation) pour observer le comportement lorsque le processeur devient le facteur limitant.

\begin{center}
\begin{tabular}{|l|r|r|r|r|r|}
\hline
\textbf{Version} & \textbf{Req/s} & \textbf{Moy.~(ms)} & \textbf{Méd.~(ms)} & \textbf{P95~(ms)} & \textbf{Max~(ms)} \\
\hline
MVC JPA         & 47,79 & 31,88 & 19,89 & 84,77 & 204 \\
WebFlux         & 47,60 & 33,76 & 22,65 & 83,96 & 217 \\
Virtual Threads & 47,84 & 31,49 & 27,37 & 50,21 &  87 \\
\hline
\end{tabular}
\end{center}

Les trois versions livrent le même débit ($\approx$~47--48~req/s)~: le facteur limitant est le
nombre de cœurs CPU, indépendamment du modèle de threading. Virtual Threads se distingue par un
P95 de 50~ms contre 84~ms pour MVC JPA et WebFlux, soit une amélioration de 40~\% en latence de
queue. Ce scénario n'implique aucun accès base de données~; la taille du pool de connexions est
sans effet sur les résultats.

\subsection{Scénario de saturation (800~VU, sans pause)}
\label{sec:saturation}

Ce scénario mesure le débit maximal soutenable sur un endpoint de lecture simple
(\texttt{GET /products/\{id\}}), sans tempo d'itération, à 800~VU simultanés.

\begin{center}
\begin{tabular}{|l|r|r|r|r|r|}
\hline
\textbf{Version} & \textbf{Req/s} & \textbf{Moy.~(ms)} & \textbf{Méd.~(ms)} & \textbf{P95~(ms)} & \textbf{Max~(ms)} \\
\hline
MVC JPA         & 16\,937 & 17,11 &  6,50 &  65,54 & 855 \\
WebFlux         & 14\,006 & 20,71 &  6,76 & 106,39 & 266 \\
Virtual Threads & 12\,299 & 23,58 &  8,71 & 103,71 & 481 \\
\hline
\end{tabular}
\end{center}

MVC JPA délivre 21~\% de débit supplémentaire par rapport à WebFlux et 38~\% par rapport à Virtual
Threads. WebFlux affiche en revanche le maximum le plus prévisible (266~ms contre 855~ms pour MVC
JPA). Ce scénario est le plus exposé à l'asymétrie des pools~; les implications sont détaillées en
section~\ref{sec:impl}.

% ══════════════════════════════════════════════════════════════════════════════
\section{Analyse et discussion}
% ══════════════════════════════════════════════════════════════════════════════

\subsection{Cohérence générale des résultats}

Les résultats sont cohérents avec la théorie des modèles de concurrence Java~:

\begin{center}
\begin{tabular}{|l|l|l|}
\hline
\textbf{Scénario} & \textbf{Résultat observé} & \textbf{Attendu par la théorie} \\
\hline
CRUD (itération cadencée) & Tous identiques en débit      & Oui — goulot dans le pacing \\
I/O-bound                 & VT $\approx$ WebFlux $\gg$ JPA & Oui — threads bloquants épuisés \\
CPU-bound                 & Tous identiques en débit      & Oui — limité par les cœurs CPU \\
Saturation                & JPA $>$ WebFlux $>$ VT        & Oui (voir §\ref{sec:impl}) \\
\hline
\end{tabular}
\end{center}

\subsection{Observations sur l'implémentation}
\label{sec:impl}

\paragraph{WebFlux — scheduler sous-optimal pour le CPU.}
Le service \texttt{CpuStressService} déporte correctement le calcul hors de la boucle
événementielle Netty via \texttt{subscribeOn}~:

\begin{lstlisting}[language=Java]
return Mono.fromCallable(() -> fibonacci(n))
           .subscribeOn(Schedulers.boundedElastic());
\end{lstlisting}

Cependant, \texttt{Schedulers.boundedElastic()} est conçu pour des opérations I/O bloquantes, non
pour du calcul CPU. Son plafond par défaut est $10 \times \text{cœurs}$ threads, ce qui provoque
une sursouscription du CPU sous charge et augmente la variance de la latence (P95 = 84~ms contre
50~ms pour Virtual Threads). Le scheduler approprié pour un travail intensif en calcul est
\texttt{Schedulers.parallel()}, borné au nombre de cœurs.

\paragraph{Asymétrie des pools de connexions et saturation.}
Les benchmarks ont été exécutés avec trois configurations de pool différentes~: 10~connexions pour
WebFlux (R2DBC par défaut), 20~connexions pour MVC JPA (HikariCP) et 50~connexions pour Virtual
Threads (HikariCP). Cette asymétrie complique l'interprétation du scénario de saturation, qui
sollicite directement la base de données.

Malgré cela, certaines conclusions restent solides. MVC JPA surpasse Virtual Threads (38~\% de
débit supplémentaire) alors que Virtual Threads dispose de 2,5~fois plus de connexions~: le goulot
de Virtual Threads n'est donc pas le pool de connexions. La cause probable est le
\emph{thread pinning}~: le driver JDBC PostgreSQL (\texttt{pgjdbc}) contient des blocs
\texttt{synchronized} dans ses chemins d'exécution, ce qui fixe (\emph{pin}) les virtual threads
à leurs threads porteurs (ForkJoinPool, borné au nombre de cœurs). Ce comportement peut être
diagnostiqué via~:

\begin{lstlisting}[language=bash]
-Djdk.tracePinnedThreads=full
\end{lstlisting}

L'avantage de MVC JPA sur WebFlux (21~\%) est quant à lui partiellement influencé par l'asymétrie
des pools (20 contre 10~connexions)~; une re-exécution avec 50~connexions pour toutes les versions
permettra de dissocier l'effet du pool de l'overhead propre au pipeline réactif.

\paragraph{Virtual Threads — comportement I/O non affecté.}
Malgré le risque de pinning, Virtual Threads performe aussi bien que WebFlux dans le scénario
I/O-bound. Le pinning n'affecte significativement que les endpoints rapides à très haute densité
de requêtes, où la durée de blocage sur le verrou \texttt{synchronized} représente une fraction
notable du temps de traitement total.

% ══════════════════════════════════════════════════════════════════════════════
\section{Synthèse comparative}
% ══════════════════════════════════════════════════════════════════════════════

\begin{center}
\begin{tabular}{|l|l|l|L{4.2cm}|}
\hline
\textbf{Scénario} & \textbf{Vainqueur} & \textbf{Classement} & \textbf{Remarque clé} \\
\hline
CRUD (200~VU)      & MVC JPA (latence)  & JPA $\approx$ WF $\approx$ VT & Max JPA instable (pool 20) \\
I/O-bound (500~VU) & WebFlux $\approx$ VT & WF $\approx$ VT $\gg$ JPA & Modèles non bloquants validés \\
CPU-bound (50~VU)  & VT (queue)         & VT $<$ JPA $\approx$ WF & Scheduler WF sous-optimal \\
Saturation (800~VU)& MVC JPA            & JPA $>$ WF $>$ VT & Pinning VT~; pools asymétriques \\
\hline
\end{tabular}
\end{center}

Aucun framework ne domine universellement. Le choix optimal dépend du profil de charge~:

\begin{itemize}
  \item \textbf{Charges I/O-bound~:} WebFlux et Virtual Threads sont équivalents. Virtual Threads
        présente l'avantage d'un code impératif sans courbe d'apprentissage réactive et sans refonte
        de la couche d'accès aux données.
  \item \textbf{Débit maximal sur endpoints rapides~:} MVC JPA reste compétitif et évite l'overhead
        des pipelines réactifs ou du scheduling des virtual threads, au prix d'un maximum de latence
        moins prévisible.
  \item \textbf{Calcul CPU intensif~:} les trois modèles livrent le même débit~; Virtual Threads
        offre la meilleure latence de queue, et WebFlux atteindrait un résultat équivalent avec
        \texttt{Schedulers.parallel()}.
\end{itemize}

% ══════════════════════════════════════════════════════════════════════════════
\section{Conclusion}
% ══════════════════════════════════════════════════════════════════════════════

Cette expérimentation produit des résultats cohérents avec la théorie des modèles de concurrence.
Spring WebFlux et Virtual Threads se comportent de façon équivalente sous charge I/O, confirmant
que Project Loom (JEP~444) atteint son objectif~: permettre d'écrire du code bloquant avec les
mêmes performances qu'un code réactif. Spring MVC + JPA demeure compétitif, voire supérieur, pour
des endpoints à forte densité de requêtes.

Deux corrections sont identifiées pour fiabiliser les résultats~:
\begin{enumerate}
  \item \textbf{Uniformiser le pool de connexions à 50} pour toutes les versions — MVC JPA
        (\texttt{hikari.maximum-pool-size: 50}), WebFlux (\texttt{spring.r2dbc.pool.max-size: 50})
        et Virtual Threads (déjà à 50) — afin d'éliminer le biais actuel du scénario de saturation.
  \item \textbf{Substituer \texttt{Schedulers.boundedElastic()} par \texttt{Schedulers.parallel()}}
        dans le service CPU de WebFlux, pour obtenir un résultat représentatif du cas d'usage
        calcul intensif.
\end{enumerate}

% ══════════════════════════════════════════════════════════════════════════════
%  Références
% ══════════════════════════════════════════════════════════════════════════════
\begin{thebibliography}{9}

\bibitem{seda}
  Welsh M., Culler D. et Brewer E.,
  \textit{SEDA~: An Architecture for Well-Conditioned, Scalable Internet Services},
  SOSP, 2001.

\bibitem{spring}
  Documentation officielle Spring Framework~: Spring MVC, Spring WebFlux et Project Reactor.
  \url{https://docs.spring.io/spring-framework/reference/}

\bibitem{r2dbc}
  Documentation Spring Data R2DBC.
  \url{https://docs.spring.io/spring-data/r2dbc/reference/}

\bibitem{jep444}
  OpenJDK, \textit{JEP~444~: Virtual Threads}, Java~21.
  \url{https://openjdk.org/jeps/444}

\bibitem{k6}
  Documentation officielle k6.
  \url{https://k6.io/docs/}

\bibitem{actuator}
  Documentation officielle Spring Boot Actuator et Micrometer.
  \url{https://docs.spring.io/spring-boot/reference/actuator/}

\bibitem{protocole}
  Bonneau M., Galice M., Jourdain G.,
  \textit{Protocole expérimental commun~: Spring MVC, Spring WebFlux et Virtual Threads},
  UQAC, 8CLD876, 28 mai 2026.

\end{thebibliography}
